\documentclass[11pt,a4paper]{article}
\usepackage[utf8]{inputenc}
\usepackage[margin=2.5cm]{geometry}
\usepackage{amsmath,amssymb}
\usepackage{booktabs}
\usepackage{graphicx}
\usepackage{hyperref}
\usepackage{xcolor}
\usepackage{float}
\usepackage{natbib}

\title{When the Model Knows It Doesn't Know:\\Activation-Guided Confidence for Educational AI}

\author{Simon McCallum\\
Norwegian University of Science and Technology (NTNU)\\
\texttt{simon.mccallum@ntnu.no}}
\date{}

\begin{document}
\maketitle

\begin{abstract}
Large language models used in educational settings are systematically overconfident: they express maximum certainty on questions they answer incorrectly, receiving large penalties under confidence-based marking (CBM) schemes that reward calibrated self-assessment.
We bridge three lines of work---pre-normalisation activation probing, asymmetric confidence scoring (HLCC), and confidence-based marking---to propose a system where an LLM's \emph{internal} uncertainty signal, extracted from pre-LayerNorm activations, is used to modulate its expressed confidence before scoring.
On a 10-question game design assessment, 7 of 11 LLMs selected maximum confidence on every question; under CBM scoring, wrong-but-confident answers incurred an average penalty of $-1.8$ points each.
Activation probing can detect when a model is unfamiliar with a question (AUC $= 0.994$ for answerability detection on Mistral-7B) and when it is likely to answer incorrectly (AUROC $= 0.70$ on MMLU).
We propose using this signal to: (1) lower the model's expressed CBM confidence level on uncertain questions, avoiding the $-2.0$ penalty; (2) trigger chain-of-thought retry on flagged questions; and (3) abstain entirely when familiarity is very low.
Preliminary results show that familiarity-gated confidence reduces Expected Calibration Error by 70\% on MMLU, and selective answering at 80\% coverage improves accuracy from 42\% to 46\%.
The combination produces an AI system that behaves more like a well-calibrated human student under CBM: confident when knowledgeable, cautious when uncertain, and willing to say ``I don't know.''
\end{abstract}

\section{Introduction}

Consider an AI teaching assistant answering student questions under confidence-based marking.
The student asks: ``In game design, what distinguishes emergent gameplay from scripted progression?''
The AI answers correctly with high confidence: $+2.0$ points.
The student then asks about an obscure mechanic from a niche genre.
The AI answers confidently and incorrectly: $-2.0$ points.

A well-calibrated human tutor would have said ``I'm not sure about that one''---selecting confidence level 1 under CBM, scoring $0.0$ instead of $-2.0$.
The AI lacks this metacognitive ability.
Its softmax confidence is equally high for both questions, because LayerNorm erases the magnitude information that carries the uncertainty signal before it reaches the output~\citep{mccallum2026activation}.

We show that this signal can be recovered and used to produce AI systems that behave more like calibrated human assessors.
The work connects three separate contributions:

\begin{enumerate}
    \item \textbf{Activation forensics}~\citep{mccallum2026activation}: Pre-normalisation activation statistics detect model uncertainty with near-perfect accuracy (AUC $= 0.994$), but the signal is erased by LayerNorm before reaching the output.

    \item \textbf{HLCC scoring}~\citep{mccallum2026hlcc}: The Hybrid Linear-Convex Confidence framework provides an asymmetric scoring rule ($R(c) = 1+c$, $P(c) = -2c^2$) that penalises overconfidence quadratically, creating a principled optimal confidence $c^* = p/[4(1-p)]$.

    \item \textbf{Confidence-based marking}: The educational CBM scheme rewards metacognitive accuracy, and our data from 129 students shows humans naturally calibrate their confidence to their knowledge ($r = 0.91$) while AI models do not.
\end{enumerate}

The bridge: if we use the activation signal to estimate $P(\text{correct})$, then compute the HLCC-optimal confidence, and discretise it into CBM levels, we get an AI system that expresses appropriate confidence.
The system is confident when the activations indicate familiarity, cautious when they do not, and abstains when the signal is very weak.

\section{Background}

\subsection{The Overconfidence Problem}

Under our 3-level CBM scheme (guessing: $+1/0$; somewhat confident: $+1.5/-0.5$; confident: $+2/-2$), 129 human students achieved a score-per-correct ratio of 1.39, while 11 LLMs achieved 1.59.
The difference is small because AI models are more accurate (83\% vs 63\%).
But the confidence \emph{pattern} is radically different: 64\% of AI models selected maximum confidence on every question, incurring the maximum $-2.0$ penalty on every wrong answer.
Humans adapted their confidence: students getting 4/10 correct used effectively level-1 confidence (score $\approx$ 4.1, close to the expected 4.0 for all-guessing), while students getting 8/10 used level 2--3 (score $\approx$ 12.7).

The HLCC framework~\citep{mccallum2026hlcc} formalises why this matters.
At 80\% accuracy, the HLCC-optimal confidence is $c^* = 0.80/[4 \times 0.20] = 1.0$---full confidence.
At 60\% accuracy, $c^* = 0.60/[4 \times 0.40] = 0.375$---cautious.
Below 50\%, the optimal strategy is near-zero confidence.
A model that always expresses $c = 1.0$ is optimal only when it is always right.

\subsection{The Activation Signal}

Pre-normalisation activation probing~\citep{mccallum2026activation} extracts five statistics from the SwiGLU feed-forward layers before LayerNorm:
\begin{enumerate}
    \item Pre-norm magnitude ($\|h\|_2$)
    \item Near-zero gate ratio
    \item Gate standard deviation
    \item SwiGLU output sparsity
    \item Excess kurtosis
\end{enumerate}

These capture \emph{structural} properties of the computation that are erased by normalisation.
On Mistral-7B, a 6-parameter logistic regression on these features achieves AUC $= 0.994$ for answerability detection (does the model recognise this question type?) and AUROC $= 0.70$ for error detection on MMLU (will the model get this question wrong?).

The signal exists inside the model but is invisible at the output.
LayerNorm divides by the standard deviation, collapsing uncertain (small magnitude) and confident (large magnitude) representations into the same unit-norm space.

\subsection{Two Signals at Different Depths}

The activation probe reveals two distinct signals~\citep{mccallum2026activation}:
\begin{itemize}
    \item \textbf{Familiarity} (layers 6--13): Does the model recognise this question type? Transfers across tasks without retraining.
    \item \textbf{Confidence} (layers 20--31): Is the model converging on a specific answer? Task-specific, does not transfer.
\end{itemize}

For CBM purposes, the familiarity signal is more useful: it tells us whether the model has been trained on this type of question, which determines whether its answer should be trusted at all.
The confidence signal tells us whether the model is committed to a specific answer, but commitment does not imply correctness.

\section{Proposed System}

\subsection{Architecture}

For each question presented to the AI system:

\begin{enumerate}
    \item Run a standard forward pass, extracting both the answer and activation statistics from all SwiGLU layers.
    \item Compute $\hat{p} = P(\text{correct})$ from the activation probe (logistic regression on 5 features, trained on non-target data).
    \item Compute HLCC-optimal confidence: $c^* = \hat{p}/[4(1-\hat{p})]$.
    \item Discretise into CBM levels:
    \begin{itemize}
        \item $c^* < 0.25$: Level 1 (guessing, $+1/0$)
        \item $0.25 \leq c^* < 0.75$: Level 2 (somewhat confident, $+1.5/-0.5$)
        \item $c^* \geq 0.75$: Level 3 (confident, $+2/-2$)
    \end{itemize}
    \item If $\hat{p} < 0.3$ (very unfamiliar), trigger one of:
    \begin{itemize}
        \item Retry with chain-of-thought prompting
        \item Abstain (select ``I don't know'')
        \item Escalate to instructor
    \end{itemize}
\end{enumerate}

The probe adds negligible computational cost---five scalar statistics from tensors already computed in the forward pass, fed to a 6-parameter linear classifier.

\subsection{Expected Behaviour}

Under this system, the AI behaves like a well-calibrated student:

\begin{table}[H]
\centering
\caption{Expected behaviour comparison under CBM scoring.}
\label{tab:behaviour}
\begin{tabular}{lccc}
\toprule
\textbf{Scenario} & \textbf{Raw AI} & \textbf{Human} & \textbf{Probe-Gated AI} \\
\midrule
Familiar + correct & Level 3 ($+2$) & Level 3 ($+2$) & Level 3 ($+2$) \\
Familiar + wrong & Level 3 ($-2$) & Level 2 ($-0.5$) & Level 2 ($-0.5$) \\
Unfamiliar + correct & Level 3 ($+2$) & Level 1 ($+1$) & Level 1 ($+1$) \\
Unfamiliar + wrong & Level 3 ($-2$) & Level 1 ($0$) & Level 1 ($0$) or abstain \\
\bottomrule
\end{tabular}
\end{table}

The critical change is in the ``unfamiliar + wrong'' cell: raw AI loses $-2.0$, the probe-gated AI loses $0.0$ (or abstains).
Over 10 questions with 2 wrong answers, this changes the CBM score from $12.0$ to $16.0$---a 33\% improvement from confidence modulation alone, with no change in accuracy.

\subsection{Connection to HLCC Theory}

The HLCC optimal confidence $c^* = p/[4(1-p)]$ has a natural interpretation in the CBM context.
The CBM level thresholds (0.25 and 0.75 in $c^*$ space) correspond to:
\begin{itemize}
    \item Level 1 ($c^* < 0.25$): $p < 0.50$---the model thinks it is more likely wrong than right.
    \item Level 2 ($0.25 \leq c^* < 0.75$): $0.50 \leq p < 0.75$---the model has partial confidence.
    \item Level 3 ($c^* \geq 0.75$): $p \geq 0.75$---the model is genuinely confident.
\end{itemize}

These thresholds are derived from the scoring theory, not arbitrary.
The transition from level 2 to level 3 at $p = 0.75$ matches the HLCC crossover point---the accuracy level at which the optimal confidence equals the accuracy itself.
Below this point, the HLCC framework recommends \emph{under-reporting} confidence; above it, confidence can equal or exceed the accuracy estimate.

\section{Preliminary Evidence}

Three pieces of evidence support the feasibility of this approach:

\textbf{1. The activation signal is strong.}
On Mistral-7B, the familiarity probe achieves AUC $= 0.994$ for answerability detection.
On MMLU, it achieves AUROC $= 0.70$ for error detection---modest, but sufficient to shift CBM levels on the most uncertain questions.

\textbf{2. Calibration improves dramatically.}
Familiarity-gated confidence reduces ECE from 0.498 to 0.152 on MMLU (70\% reduction).
This means the probe-modulated confidence actually \emph{predicts} accuracy---a basic requirement for meaningful CBM scoring.

\textbf{3. Selective answering works.}
At 80\% coverage on MMLU, selective accuracy rises from 42\% to 46\%.
At 50\% coverage, it reaches 56\%.
The probe identifies questions where the model is guessing, and these questions contribute disproportionately to errors.

\textbf{4. Humans set the benchmark.}
The human score--correct correlation of $r = 0.91$ under CBM shows what good calibration looks like.
The goal is not perfection but parity: an AI system whose CBM score distribution resembles that of a well-performing student cohort.

\section{Remaining Work}

\textbf{Full-scale MMLU evaluation.}
The probe has been tested on 100 MMLU questions locally; full 14,000-question results are in progress on A100 GPUs.
These will provide tighter estimates of selective accuracy and per-subject analysis.

\textbf{CBM simulation.}
Apply the probe-gated confidence scheme to the game design MCQs and compute the resulting CBM scores.
Compare against the raw-AI and human distributions.

\textbf{Retry experiment.}
The MMLU retry experiment (chain-of-thought on flagged questions) will quantify how much accuracy improves when the probe triggers a second attempt.
Preliminary results from TruthfulQA show that retry helps on questions where the model ``knows the domain but recalls the wrong answer.''

\textbf{70B+ scale.}
Testing on larger Mistral models (currently running on HPC) will show whether the signal scales.
Existing results (7B to 24B) show Cohen's $d$ increasing from 0.278 to 0.731.

\section{Discussion}

\subsection{The Pedagogical Argument}

An AI system that can express calibrated uncertainty is qualitatively different from one that cannot.
In educational settings---tutoring, practice assessment, study assistance---overconfident AI is actively harmful.
A student who trusts an AI that says ``I'm sure this is right'' when it is wrong develops a false sense of understanding.
A student whose AI says ``I'm not confident about this---you should check'' learns both the content and the limits of AI.

CBM provides a simple, interpretable framework for this.
Students already understand the scoring matrix: ``if you're not sure, select 'guessing' to avoid the penalty.''
Extending this to AI systems makes the calibration quality visible in a form that educators and students can directly evaluate.

\subsection{The Technical Argument}

The technical contribution is the specific pipeline: pre-norm activation statistics $\to$ logistic regression $\to$ HLCC-optimal $c^*$ $\to$ discretised CBM level.
Each step is lightweight and interpretable:
\begin{itemize}
    \item Five scalar features (not a 4096-dimensional hidden state)
    \item Six learned parameters (not a fine-tuned neural network)
    \item A closed-form optimal confidence (not a heuristic threshold)
    \item A pedagogically-motivated discretisation (not an arbitrary binning)
\end{itemize}

The system adds no additional forward passes, no model retraining, and negligible latency.
It could be deployed as a wrapper around any SwiGLU-based LLM.

\subsection{Limitations}

The probe is trained on non-target data (ARC, HellaSwag) and applied to the target (MMLU, game design).
Transfer may degrade on domains far from the training distribution.
The MMLU error detection AUROC of 0.70 is modest---the probe is better at detecting unfamiliarity than predicting specific errors.
The CBM level discretisation introduces quantisation error: the continuous HLCC confidence is more informative than a 3-level signal.

\section{Conclusion}

We propose connecting three lines of work---activation probing, HLCC scoring theory, and confidence-based marking---to create AI systems that express calibrated uncertainty in educational settings.
The activation probe detects when the model is unfamiliar with a question; the HLCC framework computes the optimal confidence; the CBM discretisation makes it legible to students and educators.
The result is an AI that behaves more like a thoughtful student: confident when it knows, cautious when it does not, and honest about the difference.

\bibliographystyle{plainnat}
\bibliography{acm-reference}

\end{document}
